\documentclass[11pt]{article}
\usepackage[T1]{fontenc}
\usepackage{lmodern,amsmath,amssymb,amsthm,booktabs,array}
\usepackage[margin=1in]{geometry}
\usepackage[hidelinks]{hyperref}
\hypersetup{pdftitle={Deterministic algorithms over represented finite fields: consequences of prime-field factorization},pdfauthor={Alec Kriebel},pdfsubject={Attributed consequence and verification record; no Zenodo deposit},pdfkeywords={finite fields, deterministic algorithms, factorization, irreducible construction}}
\newcommand{\Fp}{\mathbf F_p}
\newcommand{\Tr}{\operatorname{Tr}}
\newcommand{\PF}{\mathsf{PF}}
\newtheorem{theorem}{Theorem}
\newtheorem{proposition}[theorem]{Proposition}
\newtheorem{corollary}[theorem]{Corollary}
\newtheorem{lemma}[theorem]{Lemma}
\title{Deterministic algorithms over represented finite fields:\\consequences of prime-field factorization}
\author{Alec Kriebel\\\small ORCID: \href{https://orcid.org/0009-0001-9320-500X}{0009-0001-9320-500X}}
\date{October 6, 2026\\\small Attributed research record; no Zenodo preprint deposit}
\begin{document}
\maketitle
\begin{abstract}
We spell out uniform deterministic reductions for complete dense polynomial
factorization and prescribed-degree irreducible construction over
$K=\Fp[t]/(h)$, where $p$ is prime and $h$ is a supplied monic irreducible
polynomial. Their only oracle is prime-field polynomial factorization.
A trace-coordinate implementation and an independent prime-Frobenius
implementation handle multiplicities and inseparability. Substituting the
prime-field theorem of OpenAI gives polynomial bit complexity in the dense
input and requested output sizes. The reductions are established machinery
of Berlekamp and Shoup; the degree-multiplication construction also appears
in Rai. This note records consequences, representation details and an audit,
not a new prime-field algorithm or a priority claim. The upstream theorem
requires the general cyclotomic Hecke zero-free theorem of a separate
companion, beyond the narrower formalized quasi-Riemann theorem.
\end{abstract}

\section{Problem, provenance and exact model}
Represent $p$ in binary, put $L=\lceil\log_2p\rceil$, and supply
$h\in\Fp[T]$ monic of degree $m\ge1$. Primality of $p$ and irreducibility
of $h$ are promises, or may be checked deterministically. The field
$K=\Fp[T]/(h)$ has power basis $1,t,\ldots,t^{m-1}$; every coordinate
is a canonical residue with $L$ bits. A nonzero polynomial $f\in K[X]$
of degree $n$ is a dense list. Complete factorization returns its leading
coefficient $c$ and distinct monic irreducibles $g_i$ with positive binary
multiplicities $e_i$, such that $f=c\prod_i g_i^{e_i}$. Constants have an
empty list. Zero has no such finite unique factorization and is rejected.
The field and coefficient input uses $O((n+2)mL)$ bits; output coefficients
including the leading coefficient use $O((n+1)mL)$ and multiplicities
$O(n\log(n+1))$ bits.

For construction, a positive degree $d$ may be supplied in binary, but the
requested dense output has $(d+1)m$ coefficient slots. Polynomial in $d$
is polynomial in this output size, and is not polynomial in $\log d$.
No practical efficiency is inferred from a fixed polynomial exponent.
A dense-input convention cannot be inferred from a field cardinality alone;
the supplied defining polynomial is part of the input.

Berlekamp's reduction from general finite fields to split prime-field
polynomials is already public \cite{Berlekamp70}. Shoup's construction
reduction is public in Theorems 3.1 and 4.1 of \cite{Shoup90}.
The particular degree-$md$ lift below also occurs in Algorithm 2 and
Theorem 4.1 of Rai \cite{Rai24}. The newly supplied input is OpenAI's
prime-field theorem \cite{OAI142}. We attribute that breakthrough to its
source. Consolidating its implications does not independently solve the
base deterministic factorization problem or invent these reductions.

\section{Explicit reduction for complete factorization}
Let $P(u,L)$ be a nondecreasing bit-cost bound for a deterministic complete
factorization algorithm $\PF$ over $\Fp$ on degree at most $u$.

\begin{theorem}[Established reduction with explicit oracle accounting]
There is a uniform deterministic reduction for complete factorization
in the above representation with polynomial nonoracle overhead in $n,m,L$.
The trace version uses at most $mn$ calls to $\PF$, each of degree at most
$n$, with sum of queried degrees at most $mn^2$.
An alternative uses at most $n$ calls and sum of degrees at most $n^2$.
Neither version enumerates the field or uses an integer-factorization
or primitive-root oracle.
\end{theorem}

\subsection{Fixed algebras and trace separation}
First suppose $f$ is monic squarefree of degree $u$. If it has $s$
irreducible components, the Chinese remainder theorem gives
\[
 A=K[X]/(f)\simeq\prod_{a=1}^s\mathbf F_{p^{m u_a}},\qquad
 \sum_a u_a=u,\quad s\le u.
\]
With $q=p^m$, $q$-Frobenius is $K$-linear. Its fixed algebra
$B_q=\ker(\mathrm{Frob}_q-I)$ is $K^s$. Form its $u\times u$ matrix
and compute a kernel basis $b_1,\ldots,b_s$ by exact elimination.
For $0\le j<m$ form
\[
 c_{ij}=\sum_{k=0}^{m-1}(t^j b_i)^{p^k}\in A.
\]
Each component of $b_i$ belongs to $K$, so each component of $c_{ij}$
is $\Tr_{K/\Fp}(t^j b_i(a))\in\Fp$. This is a quotient-algebra
power computation, not a coefficientwise trace on arbitrary elements of $A$.

The trace polynomial $\sum_{k=0}^{m-1}Y^{p^k}$ is nonzero and has degree
$p^{m-1}<p^m$, so trace is nonzero, including when $p\mid m$.
For every $\delta\ne0$, multiplication by $\delta$ is bijective;
therefore $y\mapsto\Tr(\delta y)$ is nonzero. Some power-basis coordinate
$y=t^j$ detects it. Two different coordinate maps on $B_q=K^s$ differ
on some $b_i$; hence some $c_{ij}$ separates every pair of components.
Using just $\Tr(b_i)$ would not justify this claim: $\Tr(1)$ can vanish.

If $c$ has component values $a_1,\ldots,a_s\in\Fp$, its monic minimal
polynomial, over either $K$ or $\Fp$, is exactly
\[
 \mu_c(Z)=\prod_{a\in\{a_1,\ldots,a_s\}}(Z-a).
\]
A polynomial annihilates $c$ precisely when it vanishes at each coordinate;
this proves the formula, squarefreeness and degree at most $s$.
Find $\mu_c$ by the first dependence among $1,c,\ldots,c^s$ in the
$u$-dimensional $K$-space $A$. Its coefficients lie in the canonical
$\Fp$ and can be converted after checking the other basis coordinates vanish.
Factor $\mu_c$ with $\PF$. The returned factors are distinct linear
polynomials and give its occurring values. Refine each current factor block
$g$ by $\gcd(g,c-a)$ for those values $a$. These gcds partition its
component support. After all $ms$ coordinates no block contains two
components, so every block is irreducible.

\subsection{Independent prime-Frobenius variant}
The $p$-power map is $\Fp$-linear, generally not $K$-linear.
On the $um$-dimensional space with basis $t^jX^i$, form
$\mathrm{Frob}_p-I$ as an $um\times um$ matrix over $\Fp$.
Its kernel $B_p$ is $\Fp^s$, because the $p$-fixed elements of every
finite field component are exactly $\Fp$. A basis of $B_p$ separates
all coordinate pairs. For each of its $s$ basis vectors compute the
first power relation over $\Fp$ and apply the same minimal-polynomial
factorization and gcd refinements. This proves correctness independently
of the trace pairing and uses only $s$ prime calls.

For one squarefree input, the trace version uses at most $ms^3$ gcds
and the direct version at most $s^3$. A first-dependence system has
$u$ rows over $K$, or $um$ rows over $\Fp$, and at most $s+1$ columns.
A loose count for the trace version is $O(m^2u^5(L+1))$ operations in $K$.
For the direct variant one may bound the cost by
\[
 O(mu^3L+u^5)\text{ operations in }K,
 \quad O(m^3u^3+mu^4)\text{ operations in }\Fp.
\]
These counts include powering, exact elimination, dependence tracking,
root-product checks and gcds. The prime calls add at most $msP(s,L)$
or $sP(s,L)$ bits respectively. Coefficients remain canonical field
coordinates; no coefficient-height growth occurs in elimination.
Schoolbook arithmetic and polynomial extended Euclid bound a $K$ operation
by $O(m^2(L+1)^3)$ bits, so these are uniform bit-polynomial reductions.
The reference program's simple repeated dependence solves fit within the
looser polynomial bounds; operation counters are not a runtime guarantee.

\subsection{Multiplicity and inseparability}
Normalize $f$ by its leading coefficient. The standard squarefree procedure
sets $C=\gcd(f,f')$, $W=f/C$, and $i=1$. While $W\ne1$, set
$Y=\gcd(W,C)$, emit $(W/Y,i)$ if nonconstant, then set
$W=Y$, $C=C/Y$, and increment $i$. If $C\ne1$ after this loop, its
remaining multiplicities are divisible by $p$; its derivative is zero.
For $C=\sum c_k X^k$ its unique $p$th root is
\[
 H=\sum_j c_{pj}^{p^{m-1}}X^j.
\]
Indeed $a^{p^m}=a$ in $K$. Recurse on $H$ and multiply the returned
multiplicities by $p$. Copying coefficients works only when $m=1$.

To check this procedure, fix a factor of multiplicity $e$. If $p\nmid e$,
it initially occurs once in $W$ and $e-1$ times in $C$, and is emitted
at step $i=e$. If $p\mid e$, it remains wholly in $C$ and passes to
recursion with multiplicity $e/p$. This proves the result by induction.
The emitted squarefree products have disjoint irreducible supports and
sum of degrees at most $n$. Recursive input degrees sum to at most $2n$.
All gcd loops have polynomial cost; coefficient-root exponents have
$O(mL)$ bits, not $p^{m-1}$ iterations. Summing the squarefree call bounds
gives exactly the bounds in the theorem. The argument covers $p=2$, $m=1$
and wholly inseparable inputs.

For checking $h$, compute successive $T^{p^j}\bmod h$, check the gcd with
$T^{p^j}-T$ for every $j\le\lfloor m/2\rfloor$, and check
$T^{p^m}=T\bmod h$. A reducible degree-$m$ polynomial has an irreducible
factor of degree at most $m/2$, so the gcds detect it, also with repeated
factors. This uses $m$ binary modular powers and no integer factorization.

\section{Prescribed-degree construction}
Shoup \cite[Theorem 3.1]{Shoup90} gives a deterministic reduction from
constructing an irreducible degree-$N$ polynomial over $\Fp$ to $\PF$,
with polynomial overhead in $N,L$. Theorem 4.1 gives the analogous
arbitrary-base implication. For clarity an alternative known lift is as follows.
Construct $H\in\Fp[X]$ irreducible of degree $N=md$, factor it over $K$
by the preceding theorem, and take the lexicographically first factor.
The $p$-Frobenius orbit of a root has size $N$. $p^m$-Frobenius acts
by adding $m$ modulo $N$, so every orbit has size
$N/\gcd(N,m)=d$. Every factor therefore has degree $d$ and multiplicity
one. This proves the lift without finding an isomorphism or primitive
root. It is the same lift as Rai's Algorithm 2 \cite{Rai24}.
For $d=1$ return $X$; for $m=1$ direct prime construction suffices.
The extension is the explicit quotient $K[X]/(g)$ for the output $g$.

The supplementary construction audit makes Shoup's oracle use explicit:
trial division factors the numeric output degree $N$; each prime-power
construction has prime-field algebra dimension at most $N^2$; oracle
polynomial degrees are at most $N^2$, and there are
$O(NL\log(N+1))$ calls. The coefficient slots always have $L$ bits.
The bound includes $p$-power Artin--Schreier towers, odd prime-power
binomials and the special $2$-power cases. It never factors $p-1$.
Trial division of $N$ is permitted by dense output size and is not an
integer-factorization oracle on arbitrary binary inputs.
The $2$-power exception cannot be omitted: over $\mathbf F_3$,
\[
 X^4+1=(X^2+X+2)(X^2+2X+2),
\]
although $-1$ is nonsquare. The supplied construction uses the quadratic
extension in that case, following the established binomial criterion.
Factoring $H$ over $K$ has input size $O(dm^2L)$, polynomial in the
original field and dense degree-$d$ output sizes. Thus the whole construction
is bit-polynomial whenever $P$ is polynomial.

\section{The upstream input, its audit and the resulting scope}
\begin{corollary}[Consequence of OpenAI's prime-field theorem]
Using \cite[Theorem 1.1]{OAI142}, complete factorization over the represented
$K$ and degree-$d$ irreducible construction are uniform deterministic
algorithms with bit complexity polynomial in the relevant dense
input/output sizes. They require no GRH, randomness, integer-factorization
oracle or primitive-root oracle.
\end{corollary}
\begin{proof}
Substitute the cited prime-field bound
$P(u,L)=O(((u+1)L)^{10^{12}})$ into the two established reductions.
Every oracle degree, number of calls and nonoracle representation is bounded
by a fixed polynomial in the dense parameters just given.
\end{proof}

This is a consequence of the cited proof, not an independent proof of
its analytic or geometric breakthrough. Family 142's Theorem 1.2 is an
auxiliary-prime reduction with numeric prime-size parameter $E$.
Its unconditional Theorem 1.1 imports family 029's Theorem 1.2
\cite{OAI029}: all finite-order Hecke $L$-functions over every cyclotomic
field containing $\mu_{12}$ are zero-free in $\Re s>1-10^{-6}$, with the
principal pole permitted. The applications use
$F=\mathbf Q(\mu_{12q})$ and the cyclic extension $F(p^{1/q})/F$.
Abelian zeta factorization and the smoothed explicit formula produce
auxiliary primes $\ell\le c B^{20000000}$, where
$B=20+(n+1)(L+1)$. A Dirichlet-only theorem does not suffice.

We inspected the actual manuscripts at the pinned repository commit
\begin{center}\small\texttt{adc7f1241b42e322a6451854ab7e4b4c146bf78a}.\end{center}
The audit covers source hashes, the required theta, reflection, Poisson and
Mellin arguments, and the algebraic norm, divisor and splitting proof. Separate automated adversarial audits found
no substantive gap in those inspected dependencies. Limited finite checks
of auxiliary tables, local Gauss/Euler identities and exact exponent
margins are supplied; they do not prove an analytic zero-free theorem.
The huge prime-field branch was not executed. Family 003's actual Lean
scope is Dirichlet and Eisenstein-field Hecke nonvanishing; it does not
cover the needed varying cyclotomic input. No relevant Lean build was
reproduced here and this follow-on is not claimed formally verified.

For optional fixed-$r$ extraction, factor the dense polynomial $X^r-a$
and return its linear-factor roots. Complexity is polynomial in numeric
$r,m,L$, hence in $m,L$ for each fixed $r$. No polynomial-in-$\log r$
claim follows. This is merely a factoring consequence; quadratic
nonresidue and square-root consequences already in the release are not
relabeled as new results.

\section{Reproduction, limitations and disclosure}
The supplement contains exact reduction pseudocode, source/dependency
ledgers, proof-node audits and standard-library Python implementations.
The bundled prime oracle enumerates $\Fp$ only in bounded small examples;
its $O(p)$ cost is not polynomial in $\log p$ and cannot substantiate
prime-field derandomization. Independent trial-division comparison checks
377 small monic inputs; a second implementation cross-checks 346 cases.
Examples cover characteristic two, nontrivial inverse coefficient Frobenius,
repeated factors, $m=1$, constants and rejected zero/oracle errors.
The manuscript makes no assertion of practical efficiency or first priority.
The priority audit found only inherited consequences and no demonstrated
genuinely new extension. A new-solution Zenodo deposit is therefore withheld.

AI tools were used extensively in research, drafting and verification.
Automated adversarial reviews are not conventional human peer review or
refereeing; this work has not undergone such review. Alec Kriebel is the
sole listed author; no affiliation or outside endorsement is asserted.
Upstream AI-produced mathematical work is attributed to OpenAI as supplied
by its manuscript-specific citation information.

\begin{thebibliography}{9}
\bibitem{Berlekamp70} E. R. Berlekamp, \emph{Factoring polynomials over large
finite fields}, Mathematics of Computation \textbf{24} (1970), 713--735.
\href{https://doi.org/10.1090/S0025-5718-1970-0276200-X}{doi:10.1090/S0025-5718-1970-0276200-X}.
\bibitem{Shoup90} V. Shoup, \emph{New algorithms for finding irreducible
polynomials over finite fields}, Mathematics of Computation \textbf{54}
(1990), 435--447; preliminary version, FOCS 1988.
\href{https://www.shoup.net/papers/detirred.pdf}{Author-hosted primary text}.
\bibitem{Rai24} S. S. Rai, \emph{Pseudo-deterministic construction of irreducible
polynomials over finite fields}, ECCC TR24-147 (October 4, 2024), arXiv:2410.04071 (October 5, 2024).
\href{https://arxiv.org/abs/2410.04071}{Versioned public record}.
\bibitem{OAI142} OpenAI, \emph{Deterministic Polynomial Factorization over
Prime Fields}, OpenAI Math Release preprint (October 4, 2026),
Manuscript-specific citation supplied by the source.
\href{https://github.com/openai/math/blob/adc7f1241b42e322a6451854ab7e4b4c146bf78a/preprints/Deterministic-Polynomial-Factorization-over-Prime-Fields-October-4-2026/Deterministic-Polynomial-Factorization-over-Prime-Fields.pdf}{Pinned manuscript}.
\bibitem{OAI029} OpenAI, \emph{Primitive roots for every admissible integer
base}, OpenAI Math Release preprint (October 4, 2026),
Manuscript-specific citation supplied by the source.
\href{https://github.com/openai/math/blob/adc7f1241b42e322a6451854ab7e4b4c146bf78a/preprints/Primitive-roots-for-every-admissible-integer-base-October-4-2026/primitive-roots-all-integer-bases.pdf}{Pinned manuscript}.
\end{thebibliography}
\end{document}
